% concatenated from main.tex
\documentclass[11pt]{amsart}
\usepackage[utf8]{inputenc}
\usepackage{amsthm}
\usepackage{amsmath}
\usepackage{amssymb}
\usepackage{hyperref}
\usepackage{amsrefs}
%\usepackage{biblatex}
%\usepackage{sectsty}
%\sectionfont{\fontsize{15}{15}\selectfont}
%\addbibresource{refs.bib}
\newtheorem{theorem}{Theorem}[section]
\newtheorem{prop}{Proposition}[section]
\newtheorem{lemma}[theorem]{Lemma}
\newtheorem{cor}[theorem]{Corollary}
%\newtheorem*{theorem*}{Theorem}
%\theoremstyle{definition}
\newtheorem{definition}[theorem]{Definition}
\newtheorem{example}[theorem]{Example}
\newtheorem{problem}[theorem]{Problem}
%\newtheorem{notation}[notation]{Notation}
%\theoremstyle{remark}
\newtheorem{remark}[theorem]{Remark}
\newtheorem{Observation}[theorem]{Observation}
\newtheorem{question}[theorem]{Question}
\newtheorem{notation}[theorem]{Notations}
\newtheorem{answer}[theorem]{Answer}
\newtheorem{corollary}[theorem]{Corollary}
\newtheorem{proposition}[theorem]{Proposition}
\newtheorem{fact}[theorem]{Fact}
\numberwithin{equation}{section}

\newcommand{\bb}[1]{\mathbb{#1}}
\newcommand{\cl}[1]{\mathcal{#1}}
\begin{document}
\title[]{Nearly invariant subspaces in the real Hardy space}

\author{Arshad Khan}
\address{Department Of Mathematics\\
         School of Natural Sciences\\
        Shiv Nadar Institution Of Eminence\\
        Gautam Buddha Nagar - 201314\\
         Uttar Pradesh, India}
\email{ak954@snu.edu.in}

\author{Sneh Lata}
\address{Department Of Mathematics\\
         School of Natural Sciences\\
         Shiv Nadar Institution Of Eminence\\
         Gautam Buddha Nagar - 201314\\
         Uttar Pradesh, India}
\email{sneh.lata@snu.edu.in}

\author{ Dinesh Singh}
\address{Centre For Digital Sciences\\
      O. P. Jindal Global University\\
   Sonipat\\
        Haryana 131001, India}
\email{dineshsingh1@gmail.com}


\subjclass{Primary 47B02; Secondary 47B38, 30H10}

\keywords{Real Hardy space, nearly invariant subspaces, invariant subspaces, backward shift operator, inner function, almost invariant subspaces.}

\begin{abstract}
The objective of this article is to study nearly invariant subspaces of the backward shift operator on the real Hardy space. We also investigate nearly invariant subspaces with finite defect, and as a consequence, provide a characterization for almost invariant subspaces of the backward shift operator.
\end{abstract}

\maketitle

\section{Introduction} \label{intro}
Subspaces exhibiting invariance under linear operators are of central importance. However, numerous naturally arising subspaces are not completely invariant and demonstrate certain weaker forms of invariance. These weaker notions have prompted several generalizations; such as nearly invariant subspaces and almost invariant subspaces. The study of nearly invariant subspaces has emerged as a vibrant and active area of research, attracting considerable interest from the mathematical community in recent years. Nearly invariant subspaces of the Hardy space $H^2(\mathbb{D})$ were first introduced and studied by Hitt \cite{Hit} in 1988, while characterizing shift invariant subspaces of the Hardy space over an annulus. A subspace $\cl M$ of the Hardy space $H^2(\bb D)$ is called nearly invariant under the backward shift operator $T^*_z$ if $T^*_z(f)$ belongs to 
$\cl M$ whenever $f\in \cl M$ and $f(0)=0$. In Hitt's formulation, these subspaces were referred to as \textit{weakly invariant subspaces}.  The terminology \textit{nearly invariant subspaces} was later popularized by Sarason  \cite{Sar}, who in 1988 provided a new proof of Hitt's theorem by utilizing ideas based on de Branges-Rovnyak spaces \cite{DBR}.


\begin{theorem}[Hitt's theorem] \cite[Theorem 1]{Sar} \label{ThH} Suppose $\mathcal{M}$ is a non-trivial subspace of $H^2(\mathbb{D})$ which is nearly invariant under the backward shift operator $T^*_z$, and let $g$ be the function of unit norm in $\mathcal{M}$ which is orthogonal to $\mathcal{M}\cap zH^2(\mathbb{D})$ and positive at the origin. Then there exists a $T^*_z$-invariant subspace $\mathcal{N}$ of $H^2(\mathbb{D})$ such that $\mathcal{M}=g\mathcal{N}$ and $\|gf\|=\|f\|$ for all $f \in \mathcal{N}$.
\end{theorem}

The well-known examples of nearly $T^*_z$-invariant subspaces are kernel subspaces of Toeplitz operators on $H^2(\mathbb{D})$. Such kernel subspaces were also studied and described independently by Hayashi \cite{Hay} in 1986. Since the time \cite{Hay}, \cite{Hit}, and particularly \cite{Sar} appeared, nearly invariant subspaces have established themselves as an important area of research, and they can be deemed to be a proper generalization of the concept of invariant subspaces. For more insights into 
nearly invariant subspaces, one can refer to \cite{CCP}, \cite{CGP1}, \cite{CD}, \cite{CDP1}, \cite{DS}, \cite{Era}, \cite{KLD}, \cite{KLS1}, \cite{LP}, and the references therein. 

 In 2020, Chalendar, Gallardo-Guti$\rm{\acute{e}}$rrez, and Partington \cite{CGP1} introduced and studied the concept of nearly invariant subspaces with finite defect of the backward shift operator on $H^2(\mathbb{D})$. So far, lots of studies have been done on nearly invariant subspaces in Hardy spaces and related Hilbert spaces, but there appears to be no prior work addressing the study of nearly invariant subspaces in real Hardy spaces. Motivated by this gap in research, the present paper undertakes a systematic study and investigation of nearly invariant subspaces in the real Hardy space. 

 The real Hardy space $H^2_\mathbb{R}(\mathbb{D})$ consists of functions in $H^2(\mathbb{D})$ with real Fourier coefficients. It is a Hilbert space over the field $\mathbb{R}$. If $H^2(\mathbb{D})$ is viewed as a Hilbert space over the field $\mathbb{R}$, then $H^2_\mathbb{R}(\mathbb{D})$ is a closed subspace of $H^2(\mathbb{D})$. The study of real Hardy and Lebesgue spaces has drawn interest in the past few decades. For instance, in 1987, Mehanna and Singh \cite{MS} investigated invariant subspaces of the real $H^1(\mathbb{D})$ space. In 2006, Jupiter and Redett \cite{JR} studied shift invariant subspaces in the real $L^1(\mathbb{T})$ space.  In a broader context, Raghupati and Singh \cite{RS} characterized shift invariant subspaces of real $L^p(\mathbb{T})$ spaces in 2011. In 2007, Mikkola and Sasane \cite{Tol} proved Tolokonnikov’s Lemma and the inner-outer factorization in the algebra of bounded analytic functions (possibly operator-valued) on the open unit disc all of whose matrix-entries (with respect to some fixed orthonormal basis) are functions having real Fourier coefficients. A foundational result in the study of real function algebras was established by Wick \cite{BDW} in 2009, where he investigated stabilization in the real disc algebra $A_\mathbb{R}(\mathbb{D})$. He extended this result in 2010, where he addressed the problem in $H_\mathbb{R}^\infty (\mathbb{D})$ (see \cite{BDW1}).
 
 


 The organization of the paper is as follows. Section \ref{PRE} contains some essential terminologies and definitions that are required for the understanding of the paper.  It also includes some technical lemmas that will be used throughout the paper (Lemma \ref{Lem1}, Lemma  \ref{Lem2}, and Lemma \ref{Lem3}). In Section \ref{NIS},  we investigate nearly invariant subspaces in the real Hardy space (Theorem \ref{ThM1}). Section \ref{NISF} focuses on nearly invariant subspaces with finite defect of the backward shift operator on the real Hardy space (Theorem \ref{MT2} and Theorem \ref{MT3}). As a consequence, we provide a classification for almost invariant subspaces of the backward shift operator on the real Hardy space (Corollary \ref{CorA}).


\section{Preliminaries} \label{PRE}
Suppose $\mathbb{D}$ denotes the open unit disc in the complex plane $\mathbb{C}$. Recall that, the $\mathbb{C}^m$-valued Hardy space, denoted by $H^2(\mathbb{D},\mathbb{C}^m)$, is defined as 
\begin{eqnarray*} 
 H^2(\mathbb{D},\mathbb{C}^m):=\Bigg\{F(z)=\sum_{n=0}^{\infty}A_nz^n: A_n \in \mathbb{C}^m, \sum_{n=0}^{\infty}\|A_n\|_{\mathbb{C}^m}^2 < \infty \Bigg\}.   
\end{eqnarray*}
It is a Hilbert space over the field $\mathbb{C}$ with respect to the inner product defined by
\begin{eqnarray}\label{e2.1}
\langle{F,H}\rangle=\sum_{n=0}^{\infty}\langle{A_n, B_n}\rangle_{\mathbb{C}^m},
\end{eqnarray}  
where
$F(z)=\sum_{n=0}^{\infty}A_n z^n $ and $H(z)=
 \sum_{n=0}^{\infty}B_n z^n$ are in $H^2(\mathbb{D},\mathbb{C}^m)$, and $\langle \cdot,\cdot \rangle_{\mathbb{C}^m}$  denotes the standard inner product on $\mathbb{C}^m$. For $m=1$, we denote $H^2(\mathbb{D}, \mathbb{C})$ simply by $H^2(\mathbb{D})$.
 
 \begin{definition}
   The $\mathbb{C}^m$-valued real Hardy space, denoted by $H^2_{\mathbb{R}}(\mathbb{D},\mathbb{C}^m)$, consists of all functions in $H^2(\mathbb{D},\mathbb{C}^m)$ having Fourier coefficients from $\mathbb{R}^m$ in their power series representations. That is,
\begin{eqnarray*} 
 H^2_{\mathbb{R}}(\mathbb{D},\mathbb{C}^m):=\Bigg\{F(z)=\sum_{n=0}^{\infty}A_nz^n \in H^2(\mathbb{D},\mathbb{C}^m) : A_n \in \mathbb{R}^m\Bigg\}.   
\end{eqnarray*}
It is a Hilbert space over the field $\mathbb{R}$ with respect to the inner product defined in (\ref{e2.1}), where
$F(z)=\sum_{n=0}^{\infty}A_n z^n $ and $H(z)=
 \sum_{n=0}^{\infty}B_n z^n$ are in $H^2_\mathbb{R}(\mathbb{D},\mathbb{C}^m)$. 
\end{definition}
 
The real Hardy space $H^2_\mathbb{R}(\bb D,\bb C)$ 
is denoted simply as $H^2_\mathbb{R}(\bb D)$. 
If $H^2(\mathbb{D},\mathbb{C}^m)$ is viewed as a vector space over the field $\mathbb{R}$, then it is a Hilbert space with respect to the inner product defined in (\ref{e2.1}).

\begin{definition}
 The \textit{shift operator} on $H^2_{\mathbb{R}}(\mathbb{D},\mathbb{C}^m)$ , denoted by $T_z$, is defined as 
\begin{eqnarray*}
    (T_zF)(z)=zF(z), \text{ }F\in H^2_{\mathbb{R}}(\mathbb{D},\mathbb{C}^m),
\end{eqnarray*}
and its adjoint, denoted by $T^*_z$, called the \textit{backward shift operator}, is given by
\begin{eqnarray*}
    (T^*_zF)(z)=\frac{F(z)-F(0)}{z}, \text{ } F\in H^2_{\mathbb{R}}(\mathbb{D},\mathbb{C}^m).
\end{eqnarray*}   
\end{definition}
 
Throughout this article, we use the term \textit{subspace }for a closed subspace. Suppose $\mathcal{M}$ is a vector space over the field $\mathbb{R}$, then we define a set $\mathcal{M}_\mathbb{C}$ as
\begin{eqnarray*}
 \mathcal{M}_\mathbb{C}:=\big\{f + ih : f,h \in \mathcal{M}\big\}.  
\end{eqnarray*}
 It is a vector space over the field $\mathbb{C}$, called the \textit{complexification of $\mathcal{M}$}. For a vector $A=\begin{pmatrix} a_1 \\ \vdots \\ a_m\end{pmatrix}$ in $\mathbb{C}^m$, we define $\overline{A}:=\begin{pmatrix} \overline{a_1} \\ \vdots \\ \overline{a_m}\end{pmatrix}$.
  For each $F$ in $H^2(\mathbb{D},\mathbb{C}^m)$, we define a function $\widehat{F}: \mathbb{D} \longrightarrow \mathbb{C}^m$ by
\begin{eqnarray*}
    \widehat{F} (z):=\overline{F(\overline{z})}.
\end{eqnarray*}

If $F(z)=\sum_{n=0}^{\infty}A_nz^n \in H^2(\mathbb{D},\mathbb{C}^m)$, then
$
\widehat{F}(z)=\sum_{n=0}^{\infty}\overline{A_n}z^n$. It is easy to see that $\widehat{F} \in H^2(\mathbb{D},\mathbb{C}^m)$.

Given a subspace $\mathcal{N}$ of $H^2(\mathbb{D},\mathbb{C}^m)$, we define $\widehat{\mathcal{N}}:=\{\widehat{F} : F \in \mathcal{N}\}$. It is straightforward to check that $\widehat{N}$ is also a subspace of $H^2(\mathbb{D},\mathbb{C}^m)$. The Hardy spaces $H^2(\mathbb{D}, \mathbb{C}^{m})$ and $H^2(\mathbb{D}) \oplus \cdots \oplus H^2(\mathbb{D})$ ($m$ times), the orthogonal sum of $m$ copies of $H^2(\mathbb{D})$, are unitarily equivalent. The unitary is given by the map 
\begin{center}
    

$\begin{pmatrix}
a_1 \\ 
a_2\\ 
\vdots \\
a_m
\end{pmatrix}z^n
\mapsto 
\begin{pmatrix}
 a_1 z^n\\ 
a_2 z^n\\
\vdots \\
a_m z^n
\end{pmatrix},$
\end{center}
where $a_1, \dots , a_m$ are complex scalars.
Thus, each $F\in H^2(\mathbb{D},\mathbb{C}^m)$ can be identified as an $m$-tuple $(f_1, \dots , f_m)$, where $f_1, \dots , f_m$ are in $H^2(\mathbb{D})$. Also, $\widehat{F}$ is identified as $(\widehat{f_1}, \dots , \widehat{f_m})$.




\begin{definition}
The function $\phi : H^2(\mathbb{D},\mathbb{C}^m) \longrightarrow H^2(\mathbb{D},\mathbb{C}^m)$ defined by 
\begin{eqnarray*}
    \phi(F)=\frac{F + \widehat{F}}{2}, \text{ } F\in H^2(\mathbb{D},\mathbb{C}^m), 
\end{eqnarray*} 
is called symmetrisation.
\end{definition}

We observe that $\phi$ maps $H^2(\mathbb{D},\mathbb{C}^m)$ onto $H^2_{\mathbb{R}}(\mathbb{D},\mathbb{C}^m)$. The map $\phi$ is a contraction and additive, but not linear. Every function $F$ in $H^2(\mathbb{D},\mathbb{C}^m)$ can be uniquely expressed as $F=F_1 + iF_2$, where $F_1$, $F_2$ are in  $H^2_{\mathbb{R}}(\mathbb{D},\mathbb{C}^m)$. This is easily understood from the following: 
\begin{eqnarray*}
    F=\frac{F + \widehat{F}}{2} + i\bigg(\frac{F - \widehat{F}}{2i}\bigg)=\phi(F) + (F - \phi(F)).
\end{eqnarray*}
For each $F=(f_1, \dots , f_m) \in H^2(\mathbb{D},\mathbb{C}^m)$, $\phi(F)$ equals to $\big(\phi(f_1), \dots , \phi(f_m)\big)$. For a set $\mathcal{K}\subset  H^2(\mathbb{D},\mathbb{C}^m)$, we define $\phi(\mathcal{K}):=\{\phi(F): F \in \mathcal{K}\}$.
We end this section by providing some essential facts, which will be useful in the subsequent sections.

\begin{lemma}\label{Lem1}
Suppose $f$ and $h$ are two functions in $H^2(\mathbb{D})$. Then $$\langle f, h \rangle = \langle \widehat{h}, \widehat{f} \rangle.$$   
\end{lemma}
\begin{proof}
Suppose $f(z)=\sum_{n=0}^{\infty}a_{n}z^n$ and $h(z)=\sum_{n=0}^{\infty}b_{n}z^n$ are in $H^2(\mathbb{D})$. Then $$\langle f, h \rangle =\sum_{n=0}^{\infty} \langle a_n, b_n \rangle_{\mathbb{C}}=\sum_{n=0}^{\infty}\langle\overline{b_n},\text{ }\overline{a_n}\rangle_{\mathbb{C}}=\langle \widehat{h}, \widehat{f} \rangle.$$
\end{proof}

\begin{lemma}\label{Lem2}
Suppose $\psi$ is an inner function in $H^2(\mathbb{D})$, then $$\widehat{H^2(\mathbb{D})\ominus \psi
H^2(\mathbb{D})}=H^2(\mathbb{D}) \ominus \widehat{\psi} H^2(\mathbb{D}).$$
\end{lemma}
\begin{proof}
First, we observe that $\psi$ is inner implies $\widehat{\psi}$ is inner.
Now, suppose $f\in \widehat{H^2(\mathbb{D})\ominus \psi
H^2(\mathbb{D})}$, then $f=\widehat{h}$ for some $h \in H^2(\mathbb{D})\ominus \psi H^2(\mathbb{D})$. Then for any $k \in H^2(\mathbb{D})$, we have $\widehat{f}\in H^2(\mathbb{D})$ and
\begin{eqnarray*}
\langle f, \widehat{\psi} k \rangle=
\langle \widehat{h}, \widehat{\psi} k \rangle= \langle \psi \widehat{k}, h \rangle=0 
\end{eqnarray*}
This shows that $\widehat{H^2(\mathbb{D})\ominus \psi
H^2(\mathbb{D})} \subset H^2(\mathbb{D}) \ominus \widehat{\psi} H^2(\mathbb{D})$.


 Conversely, suppose $f \in H^2(\mathbb{D}) \ominus \widehat{\psi} H^2(\mathbb{D})$, then
\begin{eqnarray*}
\langle \widehat{f}, \psi h \rangle=\langle \widehat{\psi} \widehat{h}, f \rangle=0    \text{  }   \text{ for any } h \in H^2(\mathbb{D}),  
\end{eqnarray*}
which shows that $f\in \widehat{H^2(\mathbb{D})\ominus \psi H^2(\mathbb{D})}$.

Thus $$H^2(\mathbb{D}) \ominus \widehat{\psi} H^2(\mathbb{D})\subset \widehat{H^2(\mathbb{D})\ominus \psi H^2(\mathbb{D})}.$$ This completes the proof.
\end{proof}



\begin{lemma}\label{Lem3}
Suppose $\mathcal{N}$ is a non-zero subspace of $H^2(\mathbb{D})$. Then $$H^2_{\mathbb{R}}(\mathbb{D}) \ominus \phi(\mathcal{N})\subset\phi(H^2(\mathbb{D})\ominus \mathcal{N}).$$
Moreover, if $\widehat{\mathcal{N}}\subset \mathcal{N}$, then equality holds.
\end{lemma}
\begin{proof}
Suppose $f \in H^2_\mathbb{R}(\mathbb{D}) \ominus \phi(\mathcal{N})$. Then $$\langle f, \phi(h) \rangle =0  \text{ 
 } \text{for any } h \in \mathcal{N}.$$ 
This implies that for any $h \in \mathcal{N}$,
\begin{eqnarray*}
\langle f, h \rangle &=& -\langle f, \widehat{h} \rangle\\
&=& -\langle h, \widehat{f} \rangle \ \ \  \text{(by Lemma \ref{Lem1})}\\
&=& -\langle h, f \rangle  \ \ \  (\text{since} \ f =\widehat{f}).  
\end{eqnarray*}
Now, consider the following orthogonal decomposition
$$H^2(\mathbb{D})= \mathcal{N} \oplus (H^2(\mathbb{D})\ominus \mathcal{N}).$$
Then, $f$ can be expressed as $f=f_1 + f_2$, where $f_1 \in \mathcal{N}$ and $f_2 \in H^2(\mathbb{D})\ominus \mathcal{N}$. Then
\begin{eqnarray*}
\langle f_1, f_1 \rangle&=&\langle f_1 + f_2, f_1 \rangle\\
&=&\langle f, f_1 \rangle\\
&=&-\langle f_1, f \rangle\\
&=&-\langle f_1, f_1 \rangle ,   
\end{eqnarray*}
which shows that $\|f_1\|=0$; that is, $f_1=0$. Therefore $$f=\phi(f)=\phi(f_2),$$
which belongs to $\phi(H^2(\mathbb{D})\ominus \mathcal{N})$.

 
 Further, if $\widehat{\mathcal{N}}\subset \mathcal{N}$. Choose $f \in \phi(H^2(\mathbb{D})\ominus \mathcal{N})$. Then $f=\phi(k)$, for some $k \in H^2(\mathbb{D})\ominus \mathcal{N}$.
Now, for any $\phi(h) \in \phi(\mathcal{N})$,
\begin{eqnarray*}
\langle f, \phi(h) \rangle &=& \langle \phi(k), \phi(h) \rangle\\
&=& \frac{1}{4}\bigg(\langle k,h \rangle +  \langle \widehat{k},h \rangle +  \langle k, \widehat{h} \rangle + \langle \widehat{k},\widehat{h} \rangle \bigg)\\
&=& \frac{1}{4}\bigg(\langle k,h \rangle +  \langle \widehat{h},k \rangle +  \langle k, \widehat{h} \rangle + \langle h,k \rangle \bigg) \quad \text{(By Lemma \ref{Lem1})}\\
&= & 0  \quad\text{(since $\widehat{\mathcal{N}}\subset \mathcal{N}$)}
\end{eqnarray*}
which shows that $\phi(H^2(\mathbb{D})\ominus \mathcal{N}) \subset H^2_{\mathbb{R}}(\mathbb{D})
\ominus \phi(\mathcal{N})$. This completes the proof.
\end{proof} 


\section{Nearly invariant subspaces} \label{NIS}
In this section, we shall investigate the description of nearly invariant subspaces of the backward shift operator on the real Hardy space $H^2_\mathbb{R}(\mathbb{D})$.

\begin{definition}\label{defN}
 A subspace $\mathcal{M}$ of $H^2_{\mathbb{R}}(\mathbb{D})$ is called nearly invariant under the backward shift operator $T^*_z$ if $T^*_zf \in \mathcal{M}$, whenever $f \in \mathcal{M}$ and $f(0)=0$. In other words,
 \begin{eqnarray*}
     T^*_z(\mathcal{M}\cap zH^2_{\mathbb{R}}(\mathbb{D}))\subset \mathcal{M}.
 \end{eqnarray*}
 \end{definition}

Before we proceed to the main result of this section (Theorem \ref{ThM1}), we first describe shift invariant subspaces of the real Hardy space $H^2_\mathbb{R}(\mathbb{D})$. Although these subspaces were previously characterized by Mehanna and Singh \cite{MS}, in Proposition \ref{prop1} we recover the same characterization using a different techinique. We start by recalling the famous Beurling's invariant subspace theorem.

\begin{theorem}\label{BTH}
 Any non-zero shift invariant subspace $\mathcal{M}$ of $H^2(\mathbb{D})$ is of the form $\theta H^2(\mathbb{D})$, where $\theta$ is an inner function. Moreover, if $\theta_1$ and $\theta_2$ are two inner functions such that $\theta_1 H^2(\mathbb{D})=\theta_2 H^2(\mathbb{D})$, then there exists a unimodular complex scalar $\lambda$ such that $\theta_1=\lambda \theta_2$.    
\end{theorem}



\begin{proposition}\label{prop1}
Suppose $\mathcal{M}$ is a non-zero shift invariant subspace of $H^2_\mathbb{R}(\mathbb{D})$. Then there exists an inner function $\theta$ in $H^2_\mathbb{R}(\mathbb{D})$ such that $$\mathcal{M}=\theta H^2_\mathbb{R}(\mathbb{D}).$$
\end{proposition}
\begin{proof}
Consider the complexification of $\mathcal{M}$, $$\mathcal{M}_\mathbb{C}:=\mathcal{M}+i\mathcal{M}=\{f+ih : f,h \in \mathcal{M}\}.$$
Then, $\mathcal{M}_\mathbb{C}$ is a non-zero shift invariant subspace of $H^2(\mathbb{D})$. Thus, by Theorem \ref{BTH}, we have
\begin{eqnarray} \label{eqB}
\mathcal{M}_\mathbb{C}=\theta H^2(\mathbb{D}),   \end{eqnarray}
for some inner function $\theta$ in $H^2(\mathbb{D})$. Further, without loss of generality, we can choose $\theta$ in Equation (\ref{eqB}) such that the power series representation of $\theta$ has a non-zero real Fourier coefficient. Indeed, let $\theta=\sum_{n=0}^{\infty}a_nz^n$. Since $\theta$ is non-zero, there exists a non-negative integer $k$ such that $a_k\ne 0$. Let $a_k=e^{i\alpha}|a_k| \text{ for some } \alpha \in \mathbb{R}$. Then $\theta_1:=e^{-i\alpha}\theta$ is an inner function with a non-zero real Fourier coefficient and $\mathcal{M}_\mathbb{C}=\theta_1 H^2(\mathbb{D})$.
 
 Further, we also observe that $\mathcal{M}_\mathbb{C}=\widehat{\mathcal{M}_\mathbb{C}}$. Therefore 
\begin{eqnarray}\label{e3.01}
\theta H^2(\mathbb{D})=\widehat{\theta} H^2(\mathbb{D}).
\end{eqnarray}
Now, since $\theta$ is an inner function, therefore $\widehat{\theta}$ is an inner function. Thus, from Equation (\ref{e3.01}) and Theorem \ref{BTH}, there exists a non-zero unimodular complex scalar $\lambda$ such that $\theta=\lambda \widehat{\theta}$. Further, since $\theta$ has a non-zero real Fourier coefficient, therefore $\lambda=1$, and thus, we get $\theta=\widehat{\theta}$.
Therefore, from Equation (\ref{eqB}), we have 
$$\mathcal{M}=\phi(\mathcal{M}_\mathbb{C})=\phi(\theta H^2(\mathbb{D}))=\theta \phi(H^2(\mathbb{D}))=\theta H^2_\mathbb{R}(\mathbb{D}).$$
This completes the proof.
\end{proof}

\begin{theorem}\label{ThM1}
Suppose $\mathcal{M}$ is a non-zero subspace of $H_\mathbb{R}^2(\mathbb{D})$ which is nearly invariant under the backward shift operator $T^*_z$. Then there exists a $T^*_z$-invariant subspace $\mathcal{N}$ of $H_\mathbb{R}^2(\mathbb{D})$, and a function $g$ in $\mathcal{M}$, orthogonal to $\mathcal{M}\cap zH^2_\mathbb{R}(\mathbb{D})$, satisfying $g(0)>0$ such that $$\mathcal{M}=g\mathcal{N}.$$
Further, for all $gf \in \mathcal{M}$,  $\|gf\|=\|f\|$. 
\end{theorem}
\begin{proof}
We consider the complexification of $\mathcal{M}$, $$\mathcal{M}_\mathbb{C}:=\mathcal{M}+i\mathcal{M}=\{f+ih : f,h \in \mathcal{M}\}.$$
Then we observe that $\mathcal{M}_\mathbb{C}$ is a non-zero subspace of $H^2(\mathbb{D})$. Further,
since $\mathcal{M}$ is a nearly $T^*_z$-invariant subspace of $H^2_{\mathbb{R}}(\mathbb{D})$, therefore, $\mathcal{M}_\mathbb{C}$ is a nearly  $T^*_z$-invariant subspace of $H^2(\mathbb{D})$. Thus, by Theorem \ref{ThH}, there exists a $T^*_z$-invariant subspace $\mathcal{K} \subset H^2(\mathbb{D})$, and a function $g \in \mathcal{M}_{\mathbb{C}}$, orthogonal to $\mathcal{M}_\mathbb{C}\cap zH^2(\mathbb{D})$ and $g(0)>0$, such that $$\mathcal{M}_\mathbb{C}=g\mathcal{K}.$$
 Since, $\mathcal{M}_\mathbb{C}=\widehat{\mathcal{M}_\mathbb{C}}$ and $g \in \mathcal{M}_\mathbb{C}$, therefore $\widehat{g} \in \mathcal{M}_\mathbb{C}$. Also, for any $h=zf \in \mathcal{M}_\mathbb{C}\cap zH^2(\mathbb{D})$
 \begin{eqnarray*}
 \langle \widehat{g},h \rangle =\langle \widehat{g},zf \rangle = \langle z\widehat{f}, g \rangle = 0. &&\text{(by Lemma \ref{Lem1})}   
 \end{eqnarray*}
This shows that $\widehat{g}$ is orthogonal to  $\mathcal{M}_\mathbb{C}\cap zH^2(\mathbb{D})$. Now since, $\mathcal{M}_{\mathbb{C}}$ is a nearly $T^*_z$-invariant subspace of $H^2(\mathbb{D})$, therefore $\mathcal{M}_\mathbb{C}\ominus (\mathcal{M}_\mathbb{C}\cap zH^2(\mathbb{D})) $ is one dimensional. Therefore $\widehat{g}=\alpha g$, for some $\alpha \in \mathbb{C}$. Since $g(0)>0$, thus $\widehat{g}(0)=g(0)$ and hence $\alpha=1$. This concludes that $g=\widehat{g}$. Thus, $g\in \mathcal{M} \ominus (\mathcal{M}\cap zH^2_{\mathbb{R}}(\mathbb{D}))$. 

 Further, $\mathcal{M}_\mathbb{C}=g\mathcal{K}$ implies 
\begin{eqnarray*}
\mathcal{M}&=&\phi(\mathcal{M}_\mathbb{C})\\
&=&\phi(g\mathcal{K})\\&=&g\phi(\mathcal{K})\\&=&g\mathcal{N}, \text{ where } \mathcal{N}=\phi(\mathcal{K}).
\end{eqnarray*}
Now we show that $\mathcal{N}$ is a closed $T^*_z$-invariant subspace. Since $\mathcal{K}$ is a vector subspace of $H^2(\mathbb{D})$ and $\mathcal{N}=\phi(\mathcal{K})$, therefore it is easy to see that $\mathcal{N}$ is a vector subspace of $H^2_{\mathbb{R}}(\mathbb{D})$. Now, take a sequence $\{h_n\}_{n=1}^{\infty}$ in $\mathcal{N}$  such that $h_n \longrightarrow h$ for some $h$ in $H^2_{\mathbb{R}}(\mathbb{D})$, then $h_n=\phi(f_n)$ for $f_n \in \mathcal{K}$. Since $g=\widehat{g}$, therefore $\mathcal{K}=\widehat{\mathcal{K}}$, and thus $\phi(\mathcal{K}) \subset \mathcal{K}$. Hence $h \in \mathcal{K}$. Moreover,
\begin{eqnarray*}
\|gh_n-gh\|&=&\|g\phi(f_n)-gh\|\\ &=&\|\phi(f_n)-h\|\longrightarrow 0, \text{ as }n \rightarrow \infty. \quad \text{(by Theorem \ref{ThH} )}    
\end{eqnarray*}
Now, since $\mathcal{M}$ is closed and $gh_n\in \mathcal{M}$. Therefore $gh \in \mathcal{M}=g\mathcal{N},$
which shows that $h \in \mathcal{N}$. Thus $\mathcal{N}$ is closed.


 Next, take an element $h\in \mathcal{N}$. Then $h=\phi(f)$ for some $f\in \mathcal{K}$. Then
\begin{eqnarray*}
T^*_zh&=&
T^*_z\phi(f)\\&=&T^*_z\bigg(\frac{f + \widehat{f}}{2}\bigg)\\
&=&\frac{T^*_zf + \widehat{T^*_zf}}{2} \quad \text{(since $T^*_z\widehat{f}=\widehat{T^*_zf})$}\\
&\in &  \phi(\mathcal{K})=\mathcal{N},    
\end{eqnarray*}
which shows that $\mathcal{N}$ is $T^*_z$-invariant.


 Lastly, since $\phi(\mathcal{K})\subset \mathcal{K}$, therefore, for any $h=gf\in \mathcal{M}$, where $f \in \mathcal{N}= \phi(\mathcal{K})$,
\begin{eqnarray*}
\|h\|=\|gf\|=\|f\|. &&\text{(by Theorem \ref{ThH})}  
\end{eqnarray*}
This completes the proof.   
\end{proof}


\textnormal{We end this section with a few remarks that illuminate some of the structural properties of the spaces $\mathcal{N}$ and $\mathcal{K}$, which appeared in the above proof.}

\begin{remark}
 We show that $\mathcal{K}=\mathcal{N}_\mathbb{C}$, where $\mathcal{N}_\mathbb{C}$ is the complexification of $\mathcal{N}$. This is easily understood because $\mathcal{M}_\mathbb{C}=g\mathcal{K}$ and $\mathcal{M}=g\mathcal{N}$. Therefore, $\mathcal{M}_\mathbb{C}=\mathcal{M}+i\mathcal{M}
 =g\mathcal{N}_\mathbb{C}$,
 which implies $\mathcal{K}=\mathcal{N}_\mathbb{C}$.

\end{remark}

\begin{remark}\label{Rem3.3}
We establish the connection between these two inner functions. More specifically, we show that $\theta=\psi$.
In Theorem \ref{ThM1}, we observed that $\mathcal{K}$ and $\mathcal{N}$ are backward shift invariant subspaces of $H^2(\mathbb{D})$ and $H^2_{\mathbb{R}}(\mathbb{D})$, respectively.

Since, $\mathcal{K}$ is a $T^*_z$-invariant subspace of $H^2(\mathbb{D})$. Therefore, by Theorem \ref{BTH},
\begin{eqnarray*} \label{e3.1}
\mathcal{K}=H^2(\mathbb{D})\ominus\theta H^2(\mathbb{D}),   
\end{eqnarray*}
for some inner function $\theta$ in $H^2(\mathbb{D})$. In Equation (\ref{e3.1}), we can choose  $\theta$ such that the power series representation of $\theta$ has a non-zero real Fourier coefficient.

 Further, $\mathcal{N}$ is a $T^*_z$-invariant susbspace of $H^2_{\mathbb{R}}(\mathbb{D})$. Therefore, by Proposition \ref{prop1},
$$\mathcal{N}=H^2_{\mathbb{R}}(\mathbb{D})\ominus \psi H^2_{\mathbb{R}}(\mathbb{D}),$$
for some inner function $\psi$ in $H^2_{\mathbb{R}}(\mathbb{D})$.
Also, since $\mathcal{K}=\widehat{\mathcal{K}}$ and by using Lemma \ref{Lem2}, we have
\begin{eqnarray*}
H^2(\mathbb{D}) \ominus \theta H^2(\mathbb{D})&=&\widehat{H^2(\mathbb{D}) \ominus \theta H^2(\mathbb{D})}\\&=&H^2(\mathbb{D}) \ominus \widehat{\theta} H^2(\mathbb{D}),    
\end{eqnarray*}
which shows that $ \theta H^2(\mathbb{D})= \widehat{\theta} H^2(\mathbb{D})$. This, together with the fact that $\theta$ has a non-zero Fourier coefficient, implies that $\theta=\widehat{\theta}$.
Further, using Lemma \ref{Lem3}, we get $$H^2_{\mathbb{R}}(\mathbb{D}) \ominus \phi(\mathcal{K})=\phi(H^2(\mathbb{D})\ominus \mathcal{K}),$$
which implies $$\psi H^2_{\mathbb{R}}(\mathbb{D})=\phi(\theta H^2(\mathbb{D}))=\theta H^2_{\mathbb{R}}(\mathbb{D}).$$
Consequently, we get $\theta H^2(\mathbb{D})=\psi H^2(\mathbb{D})$.
Thus $\theta=\lambda \psi$ for some unimodular complex scalar $\lambda$. Since $\theta=\widehat{\theta}$ and $\psi=\widehat{\psi}$, therefore we get,
\begin{eqnarray*}
   \lambda \psi= \overline{\lambda}\psi,
\end{eqnarray*}
which shows that $\lambda=\overline{\lambda}$. Therefore $\lambda=1$, and hence $\theta=\psi$.

\end{remark}

\textnormal{The following remark concerns Hitts's theorem (Theorem \ref{ThH}) if $\mathcal{M}$ is a nearly $T^*_z$-invariant subspace of $H^2(\mathbb{D})$ and $\mathcal{M}=\widehat{\mathcal{M}}$. Then the unique function $g$ given by Hitt's theorem  in the description of $\mathcal{M}$ satisfies $g=\widehat{g}$.}
\begin{remark}\label{Rem3.5}
Consider a subspace $\mathcal{M}$ of $H^2(\mathbb{D})$ that is nearly $T^*_z$-invariant and satisfies $\mathcal{M}=\widehat{\mathcal{M}}$. Then we observe that $\phi(\mathcal{M})=\mathcal{M} \cap H^2_{\mathbb{R}}(\mathbb{D})$ and  $\phi(\mathcal{M})$ is a nearly $T^*_z$-invariant subspace of $H^2_{\mathbb{R}}(\mathbb{D})$. Therefore, by Theorem \ref{ThM1}, there exists $g \in \phi(\mathcal{M})$ satisfying $g=\widehat{g}$ such that $$\phi(\mathcal{M})=g\mathcal{N},$$
where $\mathcal{N}$ is a $T^*_z$-invariant subspace of $H^2_{\mathbb{R}}(\mathbb{D})$. Further, the condition $\mathcal{M}=\widehat{\mathcal{M}}$ implies that $\mathcal{M}=\phi(\mathcal{M}) + i\phi(\mathcal{M})$. Indeed, $\phi(\mathcal{M}) + i\phi(\mathcal{M}) \subset \mathcal{M}$, and for any $f \in \mathcal{M},$
\begin{eqnarray*}
 f=\frac{f + \widehat{f}}{2} + i\Bigg(\frac{f - \widehat{f}}{2i}\Bigg) \in \phi(\mathcal{M}) + i\phi(\mathcal{M}). 
\end{eqnarray*}
This concludes that
\begin{align*}
 \mathcal{M}= \phi(\mathcal{M}) + i\phi(\mathcal{M})=g\mathcal{N}_{\mathbb{C}}, 
\end{align*} 
where $\mathcal{N_{\mathbb{C}}}=\mathcal{N} + i\mathcal{N}$ is a $T^*_z$-invariant subspace of $H^2(\mathbb{D})$. This is what we required.
\end{remark}

     
\section{Nearly invariant subspaces with finite defect} \label{NISF}

\textnormal{In this section, we shall describe nearly invariant subspaces with finite defect under the backward shift operator on the real Hardy space $H^2_{\mathbb{R}}(\mathbb{D})$.  }


\begin{definition}\label{defD}
 A subspace $\mathcal{M}$ of $H^2_{\mathbb{R}}(\mathbb{D})$ is called nearly invariant with defect under the backward shift operator $T^*_z$ if there exists a finite dimensional subspace $\mathcal{F}$ (orthogonal to $\mathcal{M}$) of $H^2_{\mathbb{R}}(\mathbb{D})$ such that $T^*_zf \in \mathcal{M} \oplus \mathcal{F}$, whenever $f \in \mathcal{M}$ and $f(0)=0$. In other words, 
 \begin{eqnarray} \label{eD}
     T^*_z(\mathcal{M}\cap zH^2_{\mathbb{R}}(\mathbb{D}))\subset \mathcal{M} \oplus \mathcal{F}.
 \end{eqnarray}
A finite dimensional subspace $\mathcal{F}$ of the smallest dimension that satisfies (\ref{eD}) is called the \textit{defect space} of $\mathcal{M}$, and its dimension is called the \textit{defect} of $\mathcal{M}$. 
\end{definition}


\textnormal{Before we present the main theorem of this section (Theorem \ref{MT2}), we recall the following result due to Chalendar, Gallardo-Guti$\rm{\acute{e}}$rrez, and Partington from \cite{CGP1}.}


\begin{theorem}\cite[Theorem 3.2]{CGP1}\label{CGP}
Suppose $\mathcal{M}$ is a non-zero subspace of $H^2(\mathbb{D})$ which is nearly $T^*_z$-invariant with defect $n$, and let $\{e_i\}_{i=1}^{n}$ be an orthonormal basis of the defect space.
\begin{itemize}
    \item[(i)] If there exists a function in $\mathcal{M}$ which does not vanish at $0$, then there exists a $T^*_z\oplus T^*_z$-invariant subspace $\mathcal{K}\subset H^2(\mathbb{D}) \oplus H^2(\mathbb{D},\mathbb{C}^n)$ such that
\begin{eqnarray*}
    \mathcal{M}=\Big\{f=gk + z\sum_{i=1}^{n}k_ie_i : (k,k_1,\dots , k_n)\in \mathcal{K}\Big\},
\end{eqnarray*}
where $g$ is the function of unit norm in $\mathcal{M}$, orthogonal to $\mathcal{M}\cap zH^2(\mathbb{D})$, satisfying $g(0)>0$. Further, for all $f\in \mathcal{M}$,
\begin{eqnarray*}
\|f\|^2=\|k\|^2 +\sum_{i=1}^{n} \|k_i\|^2.
\end{eqnarray*}

\item[(ii)] If all functions in $\mathcal{M}$ vanish at $0$, then there exists a $T^*_z$-invariant subspace $\mathcal{K}\subset H^2(\mathbb{D},\mathbb{C}^n)$ such that
$$\mathcal{M}=\Big\{f=z\sum_{i=1}^{n}k_ie_i : (k_1,\dots, k_n)\in \mathcal{K}\Big\}.$$
 Further, for all $f \in \mathcal{M}$,
\begin{eqnarray*}
    \|f\|^2=\sum_{i=1}^{n}\|k_i\|^2.
\end{eqnarray*}
\end{itemize}

\end{theorem}


\textnormal{For the sake of notational clarity and simplicity, we restrict our attention to the case of defect one, focusing on the essential features of the result while omitting lengthy technical details. The extension to the general case follows in a similar manner.}

\begin{theorem}\label{MT2}
Suppose $\mathcal{M}$ is a non-zero subspace of $H^2_{\mathbb{R}}(\mathbb{D})$ which is nearly $T^*_z$-invariant with defect 1, and let the defect space be spanned by a unit vector $e$. 
\begin{itemize}
    \item[(i)] If there exists a function in $\mathcal{M}$ which does not vanish at $0$, then there exists a $T^*_z\oplus T^*_z$-invariant subspace $\mathcal{N}\subset H^2_{\mathbb{R}}(\mathbb{D}) \oplus H^2_\mathbb{R}(\mathbb{D})$ such that
\begin{eqnarray*}
    \mathcal{M}=\Big\{f=gh + zh_1e : (h, h_1)\in \mathcal{N}\Big\},
\end{eqnarray*}
where $g$ is the function of unit norm in $\mathcal{M}$, orthogonal to $\mathcal{M}\cap zH^2_{\mathbb{R}}(\mathbb{D})$, satisfying $g(0)>0$. Further, for all $f\in \mathcal{M}$,
\begin{eqnarray*}
\|f\|^2=\|h\|^2 + \|h_1\|^2.
\end{eqnarray*}
\item[(ii)] If all functions in $\mathcal{M}$ vanish at $0$, then there exists a $T^*_z$-invariant subspace $\mathcal{N}\subset H^2_{\mathbb{R}}(\mathbb{D})$ such that
$$\mathcal{M}=\Big\{f=zh_1e : h_1\in \mathcal{N}\Big\}.$$
 Further, for all $f \in \mathcal{M}$,
\begin{eqnarray*}
    \|f\|=\|h_1\|.
\end{eqnarray*}
\end{itemize}
\end{theorem}
\begin{proof}
Consider the complexification of $\mathcal{M}$, 
$$\mathcal{M}_\mathbb{C}:=\mathcal{M}+i\mathcal{M}=\{f+ih : f,h \in \mathcal{M}\}.$$
Then we observe that $\mathcal{M}_{\mathbb{C}}$ is a nearly $T^*_z$-invariant subspace of $H^2(\mathbb{D})$ with defect 1, and the defect space is spanned by the unit vector $e$. 

  Now, first assume that not all functions in $\mathcal{M}$ vanish at $0$. This implies that not all functions in $\mathcal{M}_{\mathbb{C}}$ vanish at $0$. Therefore, by Theorem \ref{CGP} for the case $n=1$, there exists a $T^*_z\oplus T^*_z$-invariant subspace $\mathcal{K}\subset H^2(\mathbb{D})\oplus H^2(\mathbb{D})$ such that
  $$\mathcal{M}_{\mathbb{C}}=\Big\{l=gk + zk_1e : (k,k_1)\in \mathcal{K}\Big\},$$
  where $g$ is the function of unit norm in $\mathcal{M}_\mathbb{C}$, orthogonal to $\mathcal{M}_\mathbb{C}\cap zH^2(\mathbb{D})$, satisfying $g(0)>0$. Since, $\mathcal{M}_\mathbb{C}=\widehat{\mathcal{M}_\mathbb{C}}$ and $g \in \mathcal{M}_\mathbb{C}$, therefore $\widehat{g} \in \mathcal{M}_\mathbb{C}$. Also, for any $l=zt \in \mathcal{M}_\mathbb{C}\cap zH^2(\mathbb{D})$,
 \begin{eqnarray*}
 \langle \widehat{g},l \rangle =\langle \widehat{g},zt \rangle = \langle z\widehat{t}, g \rangle = 0. &&\text{(By Lemma \ref{Lem1})}   
 \end{eqnarray*}
 This shows that $\widehat{g}$ is orthogonal to  $\mathcal{M}_\mathbb{C}\cap zH^2(\mathbb{D})$. Now, since $\mathcal{M}_{\mathbb{C}}$ is not contained in $zH^2(\mathbb{D})$, therefore $\mathcal{M}_\mathbb{C}\ominus (\mathcal{M}_\mathbb{C}\cap zH^2(\mathbb{D})) $ is one dimensional. Therefore $\widehat{g}=\alpha g$, for some $\alpha \in \mathbb{C}$. Since $g(0)>0$, thus $\widehat{g}(0)=g(0)$ and hence $\alpha=1$. This concludes that $g=\widehat{g}$. Thus, $g\in \mathcal{M} \ominus (\mathcal{M}\cap zH^2_{\mathbb{R}}(\mathbb{D}))$. 
Thus we have 
\begin{eqnarray*}
    \mathcal{M} &=& \phi(\mathcal{M}_\mathbb{C})\\
    &=&\{\phi(l): l \in \mathcal{M}_\mathbb{C} \}\\
    &=& \Big\{f=g\phi(k) + z\phi(k_1)e : (k,k_1) \in \mathcal{K}\Big\}\\
    &=& \Big\{f=gh + zh_1e : h=\phi(k), h_1=\phi(k_1)\Big\}.
\end{eqnarray*}
Now consider the set $\mathcal{N}\subset H^2_{\mathbb{R}}(\mathbb{D}) \oplus H^2_{\mathbb{R}}(\mathbb{D})$ defined by
\begin{eqnarray*}
    \mathcal{N}:=\{(h,h_1): f=gh+zh_1e \in \mathcal{M}\}.
\end{eqnarray*}
This shows that
$\mathcal{N}=\phi(\mathcal{K})$. Thus, for any $(k, k_1) \in \mathcal{K}$, $$l=gk + zk_1e  \in \mathcal{M}_{\mathbb{C}},$$
and since $\mathcal{M}_{\mathbb{C}}=\widehat{\mathcal{M}}_{\mathbb{C}}$, therefore $$\widehat{l}=g\widehat{k} + z\widehat{k}_1e  \in \mathcal{M}_{\mathbb{C}},$$
which implies $(\widehat{k},\widehat{k}_1)\in \mathcal{K}$.
This shows that $\widehat{\mathcal{K}}\subset \mathcal{K}$. Therefore, $$\mathcal{N}=\phi(\mathcal{K})\subset \mathcal{K}.$$ This leads us to have 
\begin{eqnarray*}
    \|f\|^2 &=& \|gh + h_1\|^2\\
    &=& \|h\|^2 + \|h_1\|^2. \quad \text{(by Theorem \ref{CGP})}\\
\end{eqnarray*}
Thus we conclude that any $f\in \mathcal{M}$ can be expressed as $$ f = gh + zh_1e,$$
together with norm equality
\begin{eqnarray*}
    \|f\|^2=\|h\|^2+\|h_1\|^2.
\end{eqnarray*}
Next, we show that $\mathcal{N}$ is a closed subspace of $H^2_{\mathbb{R}}(\mathbb{D}) \oplus H^2_{\mathbb{R}}(\mathbb{D})$. Since $\mathcal{K}$ is a vector subspace of $H^2(\mathbb{D}) \oplus H^2(\mathbb{D})$ and $\mathcal{N}=\phi(\mathcal{K})$, therefore it is easy to see that $\mathcal{N}$ is a vector subspace of $H^2_{\mathbb{R}}(\mathbb{D}) \oplus H^2_{\mathbb{R}}(\mathbb{D})$.
We show that $\mathcal{N}$ is closed. Take a Cauchy sequence $\{(h^n,h^n_1)\}_{n=1}^{\infty}$ in $\mathcal{N}$. Then $$f_n=gh^n + zh^n_1e \in \mathcal{M}$$
and   $$\|f_n\|^2=\|h^n\|^2 + \|h^n_1\|^2,$$ which is equal to $\|(h^n,h^n_1)\|^2$. This shows that the sequence $\{f_n\}$ is Cauchy and thus it converges to a point $f \in \mathcal{M}$. Therefore there exists $(h,h_1) \in \mathcal{K}$ such that $f=gh + zh_1e$. Thus 

\begin{eqnarray*}
\|(h^n,h^n_1)-(h,h_1)\|^2
&=&\|(h^n -h, h^n_1-h_1)\|^2\\
&=&\|(h^n-h)\|^2 + \|(h^n_1-h_1)\|^2\\
&=&\|f_n-f\|^2 \longrightarrow 0 \text{ as } n \longrightarrow \infty.
\end{eqnarray*}
This shows that $\mathcal{N}$ is closed in $H^2_{\mathbb{R}}
(\mathbb{D}) \oplus H_\mathbb{R}^2(\mathbb{D})$.

Lastly, we show that $\mathcal{N}$ is invariant under $T^*_z\oplus T^*_z$. Suppose $(h,h_1)\in \mathcal{N}$, then there exists $(k, k_1)\in \mathcal{K}$ such that $h=\phi(k), h_1=\phi(k_1)$. Thus 
\begin{eqnarray*}
    (T^*_z\oplus T^*_z) (h,h_1) &=& (T^*_z\oplus T^*_z)(\phi(k),\phi(k_1))\\
    &=& (T^*_z\phi(k),T^*_z\phi(k_1))\\
    &=& \bigg(T^*_z\bigg(\frac{k+\widehat{k}}{2}\bigg), T^*_z\bigg(\frac{k_1+\widehat{k_1}}{2}\bigg) \bigg)\\
    &=& \bigg(\frac{T^*_zk+T^*_z\widehat{k}}{2}, \frac{T^*_zk_1+T^*_z\widehat{k_1}}{2}\bigg) \\
    &=& \bigg(\frac{T^*_zk+\widehat{T^*_zk}}{2}, \frac{T^*_zk_1+\widehat{T^*_zk_1}}{2}\bigg) \quad \text{( since $T^*_z\widehat{k}=\widehat{T^*_zk}$, $T^*_z\widehat{k}_1=\widehat{T^*_zk_1}$)}\\
    &=& (\phi(T^*_zk), \phi(T^*_zk_1))\\
    &\in &\phi(\mathcal{K})=\mathcal{N},
\end{eqnarray*}
which is what we required.


 Further, for the other case, when all functions in $\mathcal{M}$ vanish at $0$, then this shows that all the functions in $\mathcal{M}_{\mathbb{C}}$ vanish at $0$. Therefore, by following similar techniques to those above, we get that 
 
 $$\mathcal{M}=\Big\{f=zh_1e : h_1\in \mathcal{N}\Big\},$$
where the subspace
\begin{eqnarray*}
    \mathcal{N}=\Big\{h_1: f=zh_1e\in \mathcal{M}\Big\}
\end{eqnarray*}
is $T^*_z$-invariant in $H^2_{\mathbb{R}}(\mathbb{D})$. Further, for all $f=zh_1e \in \mathcal{M}$,
\begin{eqnarray*}
    \|f\|=\|h_1\|.
\end{eqnarray*}
This completes the proof.
\end{proof}

\textnormal{Using arguments similar to those employed in the proof of Theorem \ref{MT2}, we can extend Theorem \ref{MT2} to the case of any finite defect. For the sake of completeness, we state the result below.}

\begin{theorem}\label{MT3}
Suppose $\mathcal{M}$ is a non-zero subspace of $H^2_{\mathbb{R}}(\mathbb{D})$ which is nearly $T^*_z$-invariant with defect $n$, and let $\{e_i\}_{i=1}^{n}$ be an orthonormal basis of the defect space. 
\begin{itemize}
    \item[(i)] If there exists a function in $\mathcal{M}$ which does not vanish at $0$, then there exists a $T^*_z\oplus T^*_z$-invariant subspace $\mathcal{N}$ of $H^2_{\mathbb{R}}(\mathbb{D}) \oplus H^2_{\mathbb{R}}(\mathbb{D},\mathbb{C}^n)$ such that
\begin{eqnarray*}
    \mathcal{M}=\Bigg\{f=gh + z\sum_{i=1}^{n}h_ie_i: (h, h_1, \dots, h_n) \in \mathcal{N}\Bigg\},
\end{eqnarray*}
where $g$ is the function of unit norm in $\mathcal{M}$, orthogonal to $\mathcal{M}\cap zH^2_{\mathbb{R}}(\mathbb{D})$, satisfying $g(0)>0$.
Further, for all $f \in \mathcal{M}$,
\begin{eqnarray*}
\|f\|^2=\|h\|^2 + \sum_{i=1}^{n}\|h_i\|^2.
\end{eqnarray*}
\item[(ii)] If all functions in $\mathcal{M}$ vanish at $0$, then there exists a $T^*_z$-invariant subspace $\mathcal{N}$ of $H^2_{\mathbb{R}}(\mathbb{D},\mathbb{C}^n)$ such that
$$\mathcal{M}=\Bigg\{f=z\sum_{i=1}^{n}h_ie_i :(h_1, \dots, h_n) \in \mathcal{N}\Bigg\}.$$
Further, for all $f\in \mathcal{M},$
\begin{eqnarray*}
    \|f\|^2=\sum_{i=1}^{n}\|h_i\|^2.
\end{eqnarray*}
\end{itemize}
\end{theorem}

\textnormal{Now we shall describe almost invariant subspaces of the backward shift operator on the real Hardy space:}

\begin{definition}\label{defA}
 A subspace $\mathcal{M}$ of $H^2_{\mathbb{R}}(\mathbb{D})$ is called almost invariant under the backward shift operator $T^*_z$ if there exists a finite dimensional subspace $\mathcal{F}$ (orthogonal to $\mathcal{M}$) of $H^2_{\mathbb{R}}(\mathbb{D})$ such that 
 \begin{eqnarray}\label{eA}
  T^*_z\mathcal{M} \subset \mathcal{M} \oplus \mathcal{F}.  
 \end{eqnarray}
A finite dimensional subspace $\mathcal{F}$ of the smallest dimension that satisfies (\ref{eA}) is called the \textit{defect space} of $\mathcal{M}$, and its dimension is called the \textit{defect} of $\mathcal{M}$. 
\end{definition}

\textnormal{Interestingly, almost $T^*_z$-invariant subspaces are linked with nearly $T^*_z$-invariant subspaces with finite defect in a natural way.
In the following corollary, we utilize Theorem \ref{MT3} to describe almost $T^*_z$-invariant subspaces in the real Hardy space $H^2_{\mathbb{R}}(\mathbb{D})$:}
\begin{corollary}\label{CorA}
A subspace $\mathcal{M}$ of $H^2_{\mathbb{R}}(\mathbb{D})$ is almost invariant under the backward shift operator $T^*_z$ with finite defect $n$ if and only if $\mathcal{M}$ is of the form described in either case (i) or case (ii) of Theorem \ref{MT3}, with an additional condition that $T^*_zg \in \mathcal{M} \oplus \mathcal{F}$ ($\mathcal{F}$ is the defect space) in case $(i)$, while case $(ii)$ is unchanged.  
\end{corollary}
\begin{proof}
 Suppose $\mathcal{M}$ is almost invariant under $T^*_z$ with defect $n$. Then $\mathcal{M}$ is nearly $T^*_z$-invariant with defect $n$. Thus, $\mathcal{M}$ satisfies the hypothesis of Theorem \ref{MT3}. Further, since $g \in \mathcal{M}$, therefore $T^*_zg \in \mathcal{M} \oplus \mathcal{F}$ also holds.

 
  Conversely, suppose $\mathcal{M}$ satisfies the hypothesis of Theorem \ref{MT3} together with $T^*_zg \in \mathcal{M} \oplus \mathcal{F}$. Then for any $f \in \mathcal{M}$
\begin{eqnarray*}
     f=gh + z\sum_{i=1}^{n} h_i e_i=gh(0) + g(h - h(0)) + z\sum_{i=1}^{n} h_i e_i.
\end{eqnarray*}
Set $l(z)=g(h - h(0)) + z\sum_{i=1}^{n} h_i e_i$. Then we observe that $l \in \mathcal{M}$ and since $l(0)=0$, therefore $T^*_zl \in \mathcal{M} \oplus \mathcal{F}$. This, along with the fact that $T^*_zg \in \mathcal{M} \oplus \mathcal{F}$ concludes that $T^*_zf \in\mathcal{M} \oplus \mathcal{F}$. This completes the proof.

\end{proof}


\subsection*{Acknowledgements}
We thank the Mathematical Sciences Foundation, Delhi, for the support and facilities needed to complete the present work. The first author thanks the University Grants Commission(UGC), India, for the support, and the second author thanks Shiv Nadar Institution of Eminence for partially supporting this research.



\subsection*{Statements and Declarations} 

\vspace{.2 cm}

\noindent {\bf Competing Interests}  The authors have no competing interests to declare that are relevant to the content of this article.  

\vspace{.2 cm}

\noindent {\bf Data availability statements} Data sharing is not applicable to this article as no new data were created or analyzed in this study.



%\bibliographystyle{plain}
\bibliography{refs.bib}

\end{document}
